% --- DWEETIYA PRAPATHAKA ---
\cleardoublepage
\chapter{द्वितीयः प्रपाठकः}

% --- Kandika 1 ---
\dnsub{॥ प्रथमा कण्डिका ॥}
\resetShloka

\twolineshloka{सयकारं सवं वापि अक्षरं स्वरितं भवेत्}
{न चोदात्तं पुरस्तस्य जात्यः स्वारः स उच्यते}

\twolineshloka{इउवर्णौ यदोदात्तावापद्येते यवौ क्वचित्}
{अनुदात्ते प्रत्यये तद् विद्यात् क्षैप्रस्य लक्षणम्}

\twolineshloka{एओ आभ्यामुदात्ताभ्यामकारो निहतश्च यः}
{अकारं यत्र लुम्पन्ति तमभिनिहितं विदुः}

\twolineshloka{उदात्तपूर्वं यत् किञ्चिच्छन्दसि स्वरितं भवेत्}
{एष सर्वबहुस्वारस्तैरोव्यञ्जन उच्यते}

\twolineshloka{अवग्रहात् परं यत्र स्वरितं स्यादनन्तरम्}
{तिरोविरामं तं विद्यादुदात्तो यद्यवग्रहः}

\twolineshloka{इकारं यत्र पश्येयुरिकारेणैव संयुतम्}
{उदात्तमनुदात्तेन प्रश्लिष्टं तं निबोधत}

\twolineshloka{स्वरे चेत् स्वरितं यत्र विवृता यत्र संहिता}
{एतत् पादान्तवृत्तस्य लक्षणं शास्त्रचोदितम्}

\twolineshloka{जात्यः स्वारः सजात्येन श्रुष्ट्यग्ने क्षैप्र उच्यते}
{तेऽमन्वताभिनिहितस्तैरोव्यञ्जन ऊतये}

\twolineshloka{तिरोविरामो विष्कभिते प्राश्लिष्टो हीन्द्र गिर्वणः}
{पादवृत्तः क ईं वेद स्वाराः सप्तैवमादयः}

\twolineshloka{उच्चादेकाक्षरात् पूर्वात् स्वर्यते यदिहाक्षरम्}
{स्वाराणां जात्यवर्ज्यानामेषा प्रकृतिरुच्यते}

\twolineshloka{चत्वारस्त्वादितः स्वाराः कम्पं पुष्यन्ति शास्त्रतः}
{उदात्ते वैकनीचे वा जुह्वो३ऽग्ने तत्र दर्शनम्}

\closesub

% --- Kandika 2 ---
\dnsub{॥ द्वितीया कण्डिका ॥}
\resetShloka

\twolineshloka{इकारान्ते पदे पूर्वे उकारे परतः स्थिते}
{हस्वं कम्पं विजानीयान्मेधावी नात्र संशयः}

\twolineshloka{इकारान्ते पदे चैवोकारद्वयपदे परे}
{दीर्घं कम्पं विजानीयाच्छग्ध्यूष्विति निदर्शनम्}

\twolineshloka{त्रयो दीर्घास्तु विज्ञेया ये च सन्ध्यक्षरेषु वै}
{समन्या पथ्या इन्द्राभ्या शेषा हस्वाः प्रकीर्तिताः}

\twolineshloka{अनेकानामुदात्तानामनुदात्तः परो यदि}
{शिवकं तं विजानीयादुदात्तः प्रत्ययो यदि}

\twolineshloka{यत्र द्विप्रभृतीनि स्युरुदात्तान्यक्षराणि तु}
{नीचं चोच्चं च परतस्तत्रोदात्तं विदुर्बुधाः}

\twolineshloka{न रेफे वा हकारे वा द्विर्भावो जायते क्वचित्}
{न च वर्गद्वितीयेषु न चतुर्थे कदाचन}

\twolineshloka{चतुर्थं तु तृतीयेन द्वितीयं प्रथमेन तु}
{आद्यमन्त्यं च मध्यं च स्वाक्षरेणैव पीडयेत्}

\twolineshloka{अनन्त्यश्च भवेत् पूर्वो अन्त्यश्च परतो यदि}
{तत्र मध्ये यमस्तिष्ठेत् सवर्णः पूर्ववर्णयोः}

\twolineshloka{वर्गान्त्याञ्छषसैः सार्धमन्तस्थैर्वापि संयुतान्}
{दृष्ट्वा यमा निवर्तन्ते आदेशिकमिवाध्वगाः}

\twolineshloka{तृतीयश्च चतुर्थश्च चतुर्थादि परं पदम्}
{द्वौ तृतीयौ चतुर्थश्च हकारादि परं पदम्}

\twolineshloka{अनुस्वारोपध्ममूलानां क्वचित् क्रमते परम्}
{रहपूर्वे संयुते चाप्युत्तरं क्रमतेऽक्षरम्}

\twolineshloka{संयोगे स्वरितं यत्र चोद्घाते पतने यथा}
{पूर्वाङ्गमादितः कृत्वा पराङ्गादौ निवेशयेत्}

\twolineshloka{संयोगे यत्र दृश्येत व्यञ्जनं विरते पदे}
{पूर्वाङ्गं तद् विजानीयाद् येनारम्भस्तत् पराङ्गम्}

\twolineshloka{संयोगात् तु परं स्वर्यं परं संयोगनायकम्}
{संयुक्तस्य तु वर्णस्य त्वस्वरं पूर्वमक्षरम्}

\twolineshloka{अनुस्वारः पदान्तश्च क्रमजः प्रत्यये स्वके}
{स्वरभक्तिस्तथा रेफः सर्वं पूर्वाङ्गमुच्यते}

\twolineshloka{पादादौ च पदादौ च संयोगाऽवग्रहेषु च}
{जशब्द इति विज्ञेयो योऽन्यः स य इति स्मृतः}

\twolineshloka{पदादावप्यविच्छेदे संयोगान्ते च तिष्ठताम्}
{वर्जयित्वा रहयाणां न जादेशः प्रदृश्यते}

\twolineshloka{स्वसंयुक्तो गुरुर्ज्ञेयः सानुस्वाराग्रिमः स्फुटः}
{अणुशेषोऽध्रिगोर्वापि युगलादिरविस्फुटः}

\closesub

% --- Kandika 3 ---
\dnsub{॥ तृतीया कण्डिका ॥}
\resetShloka

\twolineshloka{यदुदात्तमुदात्तं तद् यत् स्वरितं तत् पदे भवति नीचम्}
{यन् नीचं नीचमेव तद् यत् प्रचयस्थं तदपि नीचम्}

\twolineshloka{अयमग्निः सुतो मित्रमिदं वयमया वहाः}
{प्रियं दूतं घृतं चित्तमभिशब्दश्च नीचतः}

\twolineshloka{अर्केष्वेव सुतेष्वेव यज्ञेषु कलशेषु च}
{शतेषु सपवित्रेषु नीचादुच्चार्यते श्रुतिः}

\twolineshloka{हरिवरुणवरेण्येषु धारापुरुषेषु च स्वरति रेफः}
{विश्वानरे नकारश्च शेषा रस्वरिता नराः}

\twolineshloka{द्वौ वरुणौ वस्वरितावुदुत्तमन्त्वंवरुणौ}
{धाकारे चोरुधारायामुरुधारेव दोहने}

\twolineshloka{मात्रिकं वा द्विमात्रं वा स्वर्यते यदिहाक्षरम्}
{तस्यादितोऽर्धमात्रा वै शेषं तु परतो भवेत्}

\twolineshloka{अदीर्घं दीर्घवत् कुर्याद् द्विस्वरं यत् प्रयुज्यते}
{कम्पोत्स्वरिताभिगीतं हस्वकर्षणमेव च}

\twolineshloka{निमेषकाला मात्रा स्याद् विद्युत्कालेति चापरे}
{ऋस्वरा तुल्ययोगा वा मतिः स्यात् सोमशर्मणः}

\twolineshloka{समासेऽवग्रहं कुर्यात् पदं चात्रानुसंहितम्}
{यतोऽक्षरादिकरणं पदान्तं तस्य तं विदुः}

\twolineshloka{सर्वत्र पुत्रमित्रसखिशब्दा अद्रिशतक्रतुरवग्राह्याः}
{आदित्यविप्रजातवेदाश्च सत्पतिगोपतिवृत्रहा समुद्रश्च}

\twolineshloka{स्वर्युवो देवयवश्चारतिं देवतातये}
{चिकितिश्चुक्रुधं चैव नावगृह्णन्ति पण्डिताः}

\closesub

% --- Kandika 4 ---
\dnsub{॥ चतुर्थी कण्डिका ॥}
\resetShloka

\twolineshloka{विवृत्तयश्चतस्रो वै विज्ञेया इति मे मतम्}
{अक्षराणां नियोगेन तासां नामानि मे शृणु}

\twolineshloka{हस्वादिर्वत्सानुसृता वत्सानुसारिणी चाग्रे}
{पाकवत्युभयोर् हस्वा दीर्घवृद्धा पिपीलिका}

\twolineshloka{चतसृणां विवृत्तीनामन्तरं मात्रिकं भवेत्}
{अर्धमात्रिकमन्येषामन्येषामणुमात्रिकम्}

\twolineshloka{अभ्रमध्ये यथा विद्युद् दृश्यते मणिसूत्रवत्}
{एष च्छेदो विवृत्तीनां यथा वालेषु कर्तरी}

\twolineshloka{आपद्यते मकारो रेफोष्मसु प्रत्ययेष्वनुस्वारम्}
{यलवेषु परसवर्णं स्पर्शेषु चोत्तमापत्तिम्}

\twolineshloka{नकारान्ते पदे पूर्वे स्वरे च परतः स्थिते}
{आकारं रक्तमित्याहुर्नकारेण तु रज्यते}

\twolineshloka{नकारान्ते पदे पूर्वे व्यञ्जनैश्च यवोहिषु}
{अर्धमात्रा तु पूर्वस्य रज्यते त्वणुमात्रया}

\twolineshloka{नकारः स्वरसंयुक्तश्चतुर्युक्तो विधीयते}
{रेफो रङ्गश्च लोपश्च अनुस्वारोऽपि वा क्वचित्}

\twolineshloka{हृदयादुत्तिष्ठते रङ्गः कांस्येन समनिःस्वनः}
{मृदुश्चैव द्विमात्रश्च दधन्वाँ इति दर्शनम्}

\twolineshloka{यथा सौराष्ट्रिका नारी तक्रँ इत्यभिभाषते}
{एवं रङ्गः प्रयोक्तव्यो नारदस्य मतं यथा}

\closesub

% --- Kandika 5 ---
\dnsub{॥ पञ्चमी कण्डिका ॥}
\resetShloka

\twolineshloka{स्वरा गडदबाश्चैव ङणनमाः सहोष्मभिः}
{चतुर्णां पदजातीनां पदान्ता दश कीर्तिताः}

\twolineshloka{स्वर उच्चः स्वरो नीचः स्वरः स्वरित एव च}
{व्यञ्जनान्यनुवर्तन्ते यत्र तिष्ठति स स्वरः}

\twolineshloka{स्वरप्रधानं त्रैस्वर्यमाचार्याः प्रतिजानते}
{मणिवद् व्यञ्जनं विद्यात् सूत्रवच्च स्वरं विदुः}

\twolineshloka{दुर्बलस्य यथा राष्ट्रं हरते बलवान् नृपः}
{दुर्बलं व्यञ्जनं तद्वद् धरते बलवान् स्वरः}

\twolineshloka{ओभावश्च विवृत्तिश्च शषसा रेफ एव च}
{जिह्वामूलमुपध्मा च गतिरष्टविधोष्मणः}

\twolineshloka{स्वरप्रत्यया विवृत्तिः संहितायां तु या भवेत्}
{विसर्गस्तत्र मन्तव्यस्तालव्यो वात्र जायते}

\twolineshloka{सन्ध्यक्षरोपधे सन्धौ प्राप्तलुप्तौ यवौ यदि}
{व्यञ्जनाख्या विवृत्तिस्तु स्वराख्या प्रकृतिसंहिता}

\twolineshloka{ऊष्मान्तं विरते यत्र सन्धावो भवते च यत्}
{विवृत्तिर्या भवेत् तत्र स्वराख्यां तां विनिर्दिशेत्}

\twolineshloka{यद्योभावप्रसन्धानमुकारादिपरं पदम्}
{स्वरान्तं तादृशं विद्याद् यदन्यद् व्यक्तमूष्मणः}

\twolineshloka{प्रथमा उत्तमाश्चैव पदान्तेषु यदि स्थिताः}
{द्वितीयं स्थानमापन्नाः शषसाः प्रत्यया यदि}

\twolineshloka{प्रथमानूष्मसंयुक्तान् द्वितीयानिव दर्शयेत्}
{न चैनान् प्रतिजानीयाद् यथा मत्स्यः क्षुरोऽप्सरा}

\closesub

% --- Kandika 6 ---
\dnsub{॥ षष्ठी कण्डिका ॥}
\resetShloka

\twolineshloka{छन्दो मानं च वृत्तं च पादस्थानं त्रिकारणम्}
{ऋचः स्वच्छन्दवृत्तास्तु पादास्त्वक्षरमानतः}

\twolineshloka{ऋवर्णं स्वरभक्तिं च च्छन्दोमानेन निर्दिशेत्}
{प्रत्ययेन सहारैफं मिमीत स्वरभक्तिषु}

\twolineshloka{ऋवर्णे तु पृथग् रेफः प्रत्ययस्तु पृथग् भवेत्}
{विद्याल्लघुमृकारं तु संयुतो यदि नोष्मणा}

\twolineshloka{ऊष्मणैव हि संयुक्त ऋकारो यत्र पीड्यते}
{गुरुवर्णः स विज्ञेयः शृतासोऽत्रानुदर्शनम्}

\twolineshloka{ऋषभं च गृहीतं च बृहस्पतिं पृथिव्यां च}
{निर्ऋतिपञ्चमा ह्यत्र ऋकारा नात्र संशयः}

\twolineshloka{शषसहरादौ रेफः स्वरभक्तिर्जायते द्विपदा सन्धौ}
{इउवर्णाभ्यां हीना क्वचिदेकपदा क्रमवियुक्ता}

\twolineshloka{स्वरभक्तिस्तु द्विविधा ऋकारो रेफ एव च}
{स्वरोदा व्यञ्जनोदा च विहिताक्षरचिन्तकैः}

\twolineshloka{शषसेषु स्वरोदयां हकारे व्यञ्जनोदयाम्}
{शषसेषु विवृतां तु हकारे संवृतां विदुः}

\twolineshloka{स्वरभक्तिं प्रयुञ्जानस्त्रीन् दोषान् परिवर्जयेत्}
{इकारं वाप्युकारं च ग्रस्तदोषं तथैव च}

\twolineshloka{संयोगपरं छपरं विसर्जनीयं द्विमात्रिकं चैव}
{अवसानिकं च न लघु सानुस्वारं घुटान्तं च}

\closesub

% --- Kandika 7 ---
\dnsub{॥ सप्तमी कण्डिका ॥}
\resetShloka

\twolineshloka{पथ्यापादः प्रथमो द्वादशमात्रस्तथा तृतीयोऽपि}
{अष्टादशो द्वितीयः समापनः पञ्चदशमात्रः}

\onelineshloka{पथ्यालक्षणमुक्तं या त्वन्या सा स्मृता विपुला}

\twolineshloka{अक्षराणां लघु हस्वमसंयोगपरं यदि}
{तत् संयोगोत्तरं विद्याद् गुरु दीर्घाक्षरं च हि}

\twolineshloka{विवृत्तिर्यत्र दृश्येत स्वरस्यैवाग्रतः स्थिता}
{गुरुः स्वरः स विज्ञेयः क्षैप्रस्तत्र न विद्यते}

\twolineshloka{अन्तोदात्तमाद्युदात्तमुदात्तमनुदात्तं नीचस्वरितम्}
{मध्योदात्तं स्वरितं द्विरुदात्तमित्येता अष्टौ पदसंज्ञाः}

\twolineshloka{अष्टप्रकारं विज्ञेयं पदानां स्वरलक्षणम्}
{अग्निः सोमः प्र वो वीर्यं हविषा स्वर्वनस्पतिः}

\twolineshloka{अन्तमध्यमयोर्नापि उदमनुदनिपाते}
{आद्यात् स्वरितमुपसर्गे द्विर्नीचमाख्यात इति}

\twolineshloka{स्वरितात् पराणि यानि तानि धार्याक्षराणि तु}
{सर्वाणि प्रचयस्थानि उपोदात्तं निहन्यते}

\twolineshloka{प्रचयो यत्र दृश्येत तत्र हन्यात् स्वरं बुधः}
{स्वरितः केवलो यत्र मृदुं तत्र निपातयेत्}

\onelineshloka{पञ्चविधमाचार्यकं नाम मुखं न्यासः करणं प्रतिज्ञोच्चारणा}

\onelineshloka{अत्रोच्यते श्रेयः खलु वै सप्रतिज्ञोच्चारणा}

\twolineshloka{यस्य कस्य चिद् वर्णस्य करणं नोपलभ्यते}
{प्रतिज्ञा तत्र वोढव्या करणं हि तदात्मकम्}

\twolineshloka{तुम्बुरुनारदवसिष्ठविश्वावस्वादयश्च गन्धर्वाः}
{सामसु निभृतं करणं स्वरसौक्ष्म्यात् तेऽपि हि न कुर्युः}

\closesub

% --- Kandika 8 ---
\dnsub{॥ अष्टमी कण्डिका ॥}
\resetShloka

\twolineshloka{कौक्षेयाग्निं सदा रक्षेदश्नीयादशनं हितम्}
{जीर्णाहारः प्रबुद्धः सन्नुषसि ब्रह्म चिन्तयेत्}

\twolineshloka{शरद्विषुवतोऽतीतादुषस्युत्थानमिष्यते}
{यावद् वासन्तिकी रात्रिर्मध्यमा पर्युपस्थिता}

\twolineshloka{आम्रपालाशबिल्वानामपामार्गशिरीषयोः}
{वाग्यतः प्रातरुत्थाय भक्षयेद् दन्तधावनम्}

\twolineshloka{खदिरस्य कदम्बस्य करवीरकरञ्जयोः}
{सर्वे कण्टकिनः पुण्याः क्षीरिणश्च यशस्विनः}

\twolineshloka{तेनास्य करणे सौक्ष्म्यं माधुर्यं चोपजायते}
{वर्णाश्च कुरुते सम्यक् प्राचीनौदव्रजिर्यथा}

\twolineshloka{त्रिफलां लवणाख्येन भक्षयेच्छिक्षकः सदा}
{अग्निमेधाजनन्येषा स्वरवर्णकरी तथा}

\twolineshloka{कृत्वा चावश्यकान् धर्मान् जाठरं पर्युपास्य च}
{पीत्वा धूमं घृतं चैव शुचिर्भूत्वा ततो वदेत्}

\twolineshloka{मन्द्रेणोपक्रमेत् पूर्वं सर्वशाखास्वयं विधिः}
{सप्त मन्त्रानतिक्रम्य यथेष्टां वाचमुत्सृजेत्}

\threelineshloka{न तां समीरयेद् वाचं या प्राणमुपरोधयेत्}
{प्राणानामुपरोधेन वैस्वर्यं चोपजायते}
{स्वरव्यञ्जनमाधुर्यं लुप्यते नात्र संशयः}

\twolineshloka{कुतीर्थादागतं दग्धमपवर्णं च भक्षितम्}
{न तस्य पाठे मोक्षोऽस्ति पापाहेरिव किल्बिषात्}

\twolineshloka{सुतीर्थादागतं व्यक्तं स्वाम्नातं सुव्यवस्थितम्}
{सुस्वरेण सुवक्त्रेण प्रयुक्तं ब्रह्म राजते}

\twolineshloka{न करालो न लम्बोष्ठो न च सर्वानुनासिकः}
{गद्गदो बद्धजिह्वश्च प्रयोगान् वक्तुमर्हति}

\twolineshloka{एकचित्तो निरुद्धान्तस्तानगानविवर्जितः}
{स तु वर्णान् प्रयुञ्जीत दन्तोष्ठं यस्य शोभनम्}

\twolineshloka{पञ्च विद्यां न गृह्णन्ति चण्डाः स्तब्धाश्च ये नराः}
{अलसाश्चानरोगाश्च येषां च विसृतं मनः}

\twolineshloka{शनैर्विद्यां शनैरर्थानारोहेत् पर्वतं शनैः}
{शनैरध्वसु वर्तेत योजनानां परं व्रजेत्}

\twolineshloka{योजनानां सहस्राणि शनैर्यति पिपीलिका}
{अगच्छन् वैनतेयोऽपि पदमेकं न गच्छति}

\twolineshloka{न हि पाणिहता वाणी प्रयोगान् वक्तुमर्हति}
{बधिरस्येव तल्पस्था विदग्धा वामलोचना}

\twolineshloka{उपांशु त्वरितं चैव योऽधीते वित्रसन्निव}
{अपि रूपसहस्रेषु सन्देहेष्वेव वर्तते}

\twolineshloka{पुस्तकप्रत्ययाधीतं नाधीतं गुरुसन्निधौ}
{राजते न सभामध्ये जारगर्भा इव स्त्रियः}

\twolineshloka{अञ्जनस्य क्षयं दृष्ट्वा वल्मीकस्य तु सञ्चयम्}
{अवन्ध्यं दिवसं कुर्याद् दानाध्ययनकर्मसु}

\twolineshloka{यत् कीटैः पांसुभिः श्लक्ष्णैर्वल्मीकः क्रियते महान्}
{न तत्र बलसामर्थ्यमुद्योगस्तत्र कारणम्}

\twolineshloka{सहस्रगुणिता विद्या शतशः परिकीर्तिता}
{आगमिष्यति जिह्वाग्रे स्थलान्निम्नमिवोदकम्}

\twolineshloka{हयानामिव जात्यानामर्धरात्रार्धशायिनाम्}
{न हि विद्यार्थिनां निद्रा चिरं नेत्रेषु तिष्ठति}

\twolineshloka{न भोजनविलम्बी स्यान्न चाशननिबन्धनः}
{सुदूरमपि विद्यार्थी व्रजेद् गरुडहंसवत्}

\twolineshloka{योऽहेरिव गणाद् भीतः सौहित्यान्नरकादिव}
{राक्षसीभ्य इव स्त्रीभ्यः स विद्यामधिगच्छति}

\twolineshloka{न शठाः प्राप्नुवन्त्यर्थान् न क्लीबा न च मानिनः}
{न च लोकरवाद् भीता न च श्वःश्वःप्रतीक्षकाः}

\twolineshloka{यथा खनन् खनित्रेण भूतले वारि विन्दति}
{तथा गुरुगतां विद्यां शुश्रूषुरधिगच्छति}

\twolineshloka{शुश्रूषारहिता विद्या यद्यपि मेधागुणैः समुपयाति}
{वन्ध्येव यौवनवती न तस्य विद्या फलवती भवति}

\twolineshloka{गुरुशुश्रूषया विद्या पुष्कलेन धनेन वा}
{अथवा विद्यया विद्या चतुर्थेन न लभ्यते}

\twolineshloka{द्यूतं पुस्तकवाच्यं च नाटकेषु च सक्तिका}
{स्त्रियस्तन्द्रा च निद्रा च विद्याविघ्नकराणि षट्}

\twolineshloka{यथा व्याघ्री हरेत् पुत्रान् दंष्ट्राभिर्न च पीडयेत्}
{भीता पतनभेदाभ्यां तद्वद् वर्णान् प्रयोजयेत्}

\twolineshloka{एवं वर्णाः प्रयोक्तव्या नाव्यक्ता न च पीडिताः}
{सम्यग्वर्णप्रयोगेण ब्रह्मलोके महीयते}

\closesection